\clearpage
\fancyhead[L]{\small KOMET.}
\SourceLeaf{422}
\begin{center}
\Large KOMET.

\medskip
Göttingische gelehrte Anzeigen. Stück 28. Seite 273..278. 18. Februar 1819.
\end{center}

\par\indent Über die beiden Kometen, die \textsc{Pons} im November des vergangenen Jahres in Marseille entdeckt hat, berichten wir hier umso lieber, als sie in mehrfacher Hinsicht zu den bemerkenswertesten gehören. Der eine scheint in Deutschland zuerst an der Mannheimer Sternwarte beobachtet worden zu sein. Herr Professor \textsc{Nicolai} stellte dort die folgenden Beobachtungen an und teilte sie Herrn Hofrat \textsc{Gauss} in einem Brief vom 8. Januar mit:

\begin{center}
\setlength{\tabcolsep}{2.2pt}\begin{tabular}{c|ccc|ccc|ccc}
1818 & \multicolumn{3}{c|}{\shortstack{Mittlere Zeit\\in Mannheim}} & \multicolumn{3}{c|}{\shortstack{Rektaszension}} & \multicolumn{3}{c}{\shortstack{Nordliche\\Deklination}} \\
\hline
Dez. 22 & $6^{\mathrm u}$ & $51^{\mathrm m}$ & $48^{\mathrm s}$ & $326^\circ$ & $18'$ & $13''$ & $2^\circ$ & $54'$ & $5''$ \\
        & 8     & 45     & 37     & 326         & 16    & 59     & 2         & 52    & 24    \\
     23 & 7     & 11     & 25     & 326         & 3     & 39     & 2         & 40    & 26    \\
     24 & 7     & 14     & 3      & 325         & 48    & 33     & 2         & 26    & 59    \\
     25 & 6     & 48     & 49     & 325         & 32    & 49     & 2         & 13    & 34    \\
     29 & 7     & 0      & 31     & 324         & 19    & 58     & 1         & 10    & 57    \\
\end{tabular}
\end{center}

\Needspace{10\baselineskip}\par\indent Dem Schreiben waren zugleich die folgenden parabolischen Elemente beigefügt, die Herr Professor \textsc{Nicolai} vorläufig aus den Beobachtungen vom 22., 25. und 29. Dezember berechnet hatte.

\begin{center}
\begin{tabular}{l l}
Zeit der Sonnennähe im Januar 1819 . . . . . . . & 25.0000 in M. \\
Logarithmus des kleinsten Abstands . . . . . . & 0.52933 \\
Länge der Sonnennähe . . . . . . . . . . & $145^\circ$ $49'$ 10 \\
Länge des aufsteigenden Knoten . . . . . . . & 330 \quad 14 \quad 17 \\
Neigung der Bahn . . . . . . . . . . . . & 14 \quad 59 \quad 6 \\
\multicolumn{2}{l}{Richtung: \emph{rechtläufig.}} \\
\end{tabular}
\end{center}

\Needspace{14\baselineskip}\par\indent An der hiesigen Sternwarte machte Herr Professor \textsc{Harding} die folgenden Beobachtungen:

\begin{center}
\setlength{\tabcolsep}{2.2pt}\begin{tabular}{c|ccc|ccc|ccc l}
& \multicolumn{3}{c|}{\shortstack{Mittlere Zeit\\in Göttingen}} & \multicolumn{3}{c|}{\shortstack{Rektaszension}} & \multicolumn{3}{c}{Deklination} & \\
\hline
1818 Dez. 25 & $6^{\mathrm u}$ & $39^{\mathrm m}$ & $52^{\mathrm s}$ & $325^\circ$ & $33'$ & $30''$ & $2^\circ$ & $14'$ & $49''$ & Nordl. \\
          26 & 5     & 58     & 32     & 325         & 17    & 8      & 2         & 0     & 31     & \\
          27 & 6     & 7      & 10     & 324         & 59    & 15     & 1         & 45    & 5      & \\
          28 & 6     & 24     & 8      & 324         & 40    & 17     & 1         & 27    & 31     & \\
1819 Jan.  1 & 6     & 36     & 16     & 323         & 11    & 48     & 0         & 14    & 58     & \\
           7 & 6     & 39     & 41     & 319         & 53    & 59     & 2         & 19    & 59     & Südl. \\
           8 & 7     & 37     & 34     & 319         & 15    & 28     & 2         & 51    & 49     & \\
          12 & 6     & 6      & 7      & 315         & 36    & 53     & 5         & 35    & 24     & \\
\end{tabular}
\end{center}

\Needspace{11\baselineskip}\par\indent An der Seeberger Sternwarte beobachtete Herr Professor \textsc{Encke} diesen Kometen fünfmal:

\begin{center}
\setlength{\tabcolsep}{2.2pt}\begin{tabular}{c|ccc|ccc|ccc l}
& \multicolumn{3}{c|}{\shortstack{Mittlere Zeit\\in Seeberg}} & \multicolumn{3}{c|}{\shortstack{Rektaszension}} & \multicolumn{3}{c}{Deklination} & \\
\hline
1819 Jan. 1 & $6^{\mathrm u}$ & $38^{\mathrm m}$ & $46^{\mathrm s}$ & $323^\circ$ & $11'$ & $36''$ & $0^\circ$ & $14'$ & $38''$ & Nördl. \\
          4 & 7     & 26     & 56     & 321         & 43    & 43     & 0         & 54    & 59     & Südl. \\
          5 & 7     & 12     & 12     & 321         & 10    & 3      & 1         & 20    & 57     & \\
          6 & 5     & 49     & 48     & 320         & 35    & 14     & 1         & 47    & 57     & \\
         12 & 6     & 16     & 32     & 315         & 35    & 20     & 5         & 35    & 26     & \\
\end{tabular}
\end{center}

\SourceLeaf{423}
\expandafter\def\csname HistoricalNoteText423\endcsname{wahrscheinlich Seeberger Zeit, was Herr \textsc{Encke} nicht ausdrücklich angegeben hatte.}
\par\indent Die Ergebnisse, die Herr Professor \textsc{Encke} aus seinen eigenen und den Mannheimer Beobachtungen gewonnen und Herrn Hofrat \textsc{Gauss} in einem Brief vom 5. Februar mitgeteilt hat, verdienen größte Aufmerksamkeit. Obwohl diese Beobachtungen nur drei Wochen umfassen, ließen sie sich durch eine Parabel nur mit Fehlern von mehreren Minuten darstellen, die bei diesen Beobachtungen unzulässig waren. Nach mehreren Versuchen fand Herr \textsc{Encke} unter Einbeziehung der freilich nur sehr groben Angaben von \textsc{Pons} eine elliptische Bahn mit einer Umlaufzeit von 3 Jahren und 7 Monaten, die mit allen Beobachtungen ausgezeichnet übereinstimmte.

\Needspace{16\baselineskip}\par\indent Die folgende Übersicht zeigt die Unterschiede zwischen den in Mannheim und Seeberg beobachteten und den berechneten Positionen:

\begin{center}
\begin{tabular}{c|c|c|l}
& Rektaszension & Deklination & Beobachter \\
\hline
Dez. 22 & $+\ 7''6$  & $-\ 4''7$  & Nicolai \\
     23 & $+\ 0.3$   & $+17.3$    & --- \\
     24 & $-\ 1.0$   & $+11.7$    & --- \\
     25 & $+11.2$    & $-\ 5.5$   & --- \\
     29 & $+14.6$    & $-\ 1.8$   & --- \\
Jan.  1 & $+15.6$    & $+14.8$    & Encke \\
      4 & $+26.5$    & $+27.2$    & --- \\
      5 & $+30.6$    & $+12.9$    & --- \\
      6 & $+36.1$    & $+23.8$    & --- \\
     12 & $+\ 7.1$   & $+\ 3.0$   & --- \\
\end{tabular}
\end{center}

\par\indent Ohne die Angaben von \textsc{Pons} hätte sich eine noch merklich bessere Übereinstimmung erzielen lassen.

\Needspace{14\baselineskip}\par\indent Die elliptischen Elemente sind folgende:

\begin{center}
\begin{tabular}{ll}
Zeit der Sonnennähe 1819. Januar . . & $27.13417$\footnotemark \\
Länge der Sonnennähe . . . . . & $156^\circ\;14'\;8''$ \\
Länge des Knotens . . . . . . . & $334\;18\;8$ \\
\multicolumn{2}{l}{(beide auf das mittlere Äquinoktium von 1819 bezogen)} \\
\end{tabular}\par\smallskip
\begin{tabular}{ll}
Neigung der Bahn . . . . . . . & $13^\circ\;42'\;30''$ \\
Logarithmus der großen Halbachse . & $0.3697758$ \\
Exzentrizitätswinkel . . . . . . & $58^\circ\;57'\;24''$ \\
\end{tabular}
\end{center}
\footnotetext{\csname HistoricalNoteText423\endcsname}

\par\indent Hoffentlich werden in Zukunft noch frühere \emph{gute} Beobachtungen aus Frankreich oder Italien bekannt, sodass sich die Unsicherheit der Bahnbestimmung enger eingrenzen lässt. Noch wichtiger wäre es, wenn in südlichen Gegenden nach dem Durchgang durch die Sonnennähe noch einige gute Beobachtungen gemacht werden könnten. In unseren Breiten, wo sich der Komet nur noch wenig über den südlichen Horizont erheben wird, besteht dazu keine Hoffnung.

\Needspace{15\baselineskip}\par\indent Höchst bemerkenswert ist die große Übereinstimmung dieser Elemente mit denen des ersten Kometen von 1805. Zum Vergleich geben wir sie hier so an, wie sie Herr Professor \textsc{Bessel} damals für eine Parabel berechnet hat. Wir bemerken nur noch, dass schon damals der Vergleich der Beobachtungen mit der Rechnung fast unverkennbare Spuren einer Abweichung der Bahn jenes Kometen von der Parabel zeigte. Diesen Spuren scheint jedoch kein Astronom weiter nachgegangen zu sein.

\begin{center}
\begin{tabular}{l l}
\multicolumn{2}{l}{Durchgang durch die Sonnennähe 1818 November} \\
\multicolumn{2}{l}{18.13782 Pariser Zeit} \\
Länge der Sonnennähe . . . . . . . . . . . . & $147^\circ$ $51'$ $28''$ \\
Länge des aufsteigenden Kometen . . . . . . . . & 344\ \ 37\ \ 19 \\
Neigung der Bahn . . . . . . . . . . . . . . & 15\ \ 36\ \ 36 \\
Logarithmus des kleinsten Abstands . . . . . . . & 9.57820 \\
\end{tabular}
\end{center}

\SourceLeaf{424}\noindent (Mit den Elementen von \textsc{Encke} für den Kometen von 1818. 1819 ergibt sich für diesen Logarithmus 9.52579). Die Frage verdient gewiss eine sorgfältige Untersuchung, ob der jetzige Komet seit 1805 solche Störungen erfahren haben kann, dass sich dadurch die Unterschiede in den Elementen erklären lassen, und ob die beiden Kometen demnach als identisch anzusehen sind.

\par\indent Den anderen von \textsc{Pons} entdeckten Kometen fand seinerseits auch Herr Professor \textsc{Bessel} am 22. Dezember an der Königsberger Sternwarte auf.

\Needspace{17\baselineskip}\par\indent Die folgenden Beobachtungen teilte dieser Astronom Herrn Hofrat \textsc{Gauss} in einem Brief vom 14. Januar mit:

\begin{center}
\setlength{\tabcolsep}{2.2pt}\begin{tabular}{c|ccc|ccc|ccc}
& \multicolumn{3}{c|}{\shortstack{Mittlere Zeit\\in Königsberg}} & \multicolumn{3}{c|}{\shortstack{Rektaszension}} & \multicolumn{3}{c}{\shortstack{Nordliche\\Deklination}} \\
\hline
1818 Dez. 22 & $7^{\mathrm u}$ & $6^{\mathrm m}$  & $51^{\mathrm s}$ & $303^\circ$ & $1'$  & $21''7$ & $36^\circ$ & $48'$ & $20''2$ \\
             & 7     & 53     & 35     & 303         & 10    & 14.7    & 36         & 48    & 29.8    \\
             & 10    & 21     & 47     & 303         & 37    & 29.1    & 36         & 51    & $0.3\!:$ \\
          24 & 18    & 22     & 50     & 311         & 56    & $29.1\!:$ & \multicolumn{3}{c}{.\quad .\quad .\quad .} \\
          25 & 6     & 22     & 20     & 313         & 17    & 17.2    & 37         & 7     & 53.1    \\
          26 & 6     & 4      & 56     & 315         & 38    & 48.2    & 37         & 4     & 1.2     \\
             & 10    & 33     & 0      & 316         & 3     & 18.1    & 37         & 2     & 35.2    \\
          27 & 6     & 55     & 0      & 317         & 39    & 48.8    & 36         & 58    & 45.6    \\
          28 & 6     & 0      & 22     & 319         & 24    & $50.7\!:$ & 36       & 53    & $31.6\!:$ \\
1819 Jan.  1 & 11    & 11     & 4      & 324         & 39    & 12.3    & 36         & 22    & $34.2\!::$ \\
           2 & 6     & 5      & 6      & 325         & 22    & 3.4     & 36         & 15    & 54.3    \\
\end{tabular}
\end{center}

\Needspace{12\baselineskip}\par\indent Dem Brief waren die folgenden vorläufigen parabolischen Elemente beigefügt, die Herr Professor \textsc{Bessel} auf die Beobachtungen vom 22. und 27. Dezember und vom 2. Januar gestützt hatte:

\begin{center}
\begin{tabular}{l l}
\multicolumn{2}{l}{Zeit der Sonnenähe 1818 Dezember 4.0968 Pariser Zeit.} \\
Aufsteigender Knoten . . . . . . . . . . . . & $90^\circ$ $7'$ $29''$ \\
Neigung der Bahn . . . . . . . . . . . . . & 117\ \ 19\ \ 10 \\
Abstand des Perihels vom Knoten . . . . . . . . & 347\ \ \ 0\ \ 24 \\
Logarithmus des kleinsten Abstands . . . . . . . & 9.928324 \\
\end{tabular}
\end{center}

\Needspace{14\baselineskip}\par\indent Mithilfe dieser Elemente fand Herr Hofrat \textsc{Gauss} den Kometen sofort als sehr blassen Nebelfleck auf. Die Elemente ergaben eine um etwa 40 zu große Rektaszension und eine um $10'$ zu große Deklination. Daraufhin stellte Herr Professor \textsc{Harding} an der hiesigen Sternwarte noch die folgenden Beobachtungen an:

\begin{center}
\setlength{\tabcolsep}{2.2pt}\begin{tabular}{c|ccc|ccc|ccc}
& \multicolumn{3}{c|}{\shortstack{Mittlere Zeit\\in Göttingen}} & \multicolumn{3}{c|}{\shortstack{Rektaszension}} & \multicolumn{3}{c}{\shortstack{Nordliche\\Deklination}} \\
\hline
1819 Jan. 26 & $7^{\mathrm u}$ & $19^{\mathrm m}$ & $5^{\mathrm s}$  & $335^\circ$ & $23'$ & $7''$  & $35^\circ$ & $18'$ & $29''$ \\
          27 & 7     & 2      & 35     & 335         & 36    & 4      & 35         & 18    & 50     \\
          28 & 7     & 22     & 16     & 335         & 50    & 38     & 35         & 20    & 58     \\
          29 & 7     & 27     & 56     & 336         & 4     & 38     & 35         & 22    & 32     \\
          30 & 6     & 54     & 10     & 336         & 16    & 33     & 35         & 23    & 31     \\
\end{tabular}
\end{center}

\par\indent Da dieser Komet, der seit seiner Entdeckung durch seine starke geozentrische Bewegung auffällt, immer lichtschwächer wird, sind kaum noch spätere Beobachtungen zu erwarten. Ob auch bei diesem Kometen die vorhandenen Beobachtungen ausreichen, um die Abweichung der Bahn von der Parabel zu bestimmen, und zu welchen Ergebnissen das führt, wird eine künftige Untersuchung zeigen.
